\cotefile{157-3}
\scanpage{82}
\mspage{\unsure{31|34}}

Supposons p. ex. $f$ propre, considérons, \unsure{pour} $\deleted{\unsure{y}}s\in\mathrm{Ob}\,\unsure{S|\mathscr S}$
\diagram{\begin{tikzcd}[ampersand replacement=\&]
X_s \arrow[r,no head,"f_s"] \arrow[d,"i_s^X"'] \& Y_{\deleted{\unsure{y}}s} \arrow[d,"i_s^Y"]\\
X_{/\deleted{\unsure{y}}s} \arrow[r,no head,"f_{/s}"] \& Y_{/s},
\end{tikzcd}}
comme $X,Y$ propres, \deleted{\unsure{$X_s$} \illegible{1} \unsure{$Y_s$}} les flèches verticales sont dans $W$ (il suffirait de beaucoup moins que la propreté \ldots), donc $f_s\in W\Longrightarrow f_{/s}\in W$.

Ainsi, ``$f$ est dans $W$ loc. sur $S$'', ce qui implique $f\in W$ par $W(2)$.

\uline{Prop. 6} Soit
\diagram{\begin{tikzcd}[ampersand replacement=\&,column sep=small,row sep=small]
X \arrow[d,"f"'] \& X' \arrow[l,"{\unsure{g|q}}"'] \arrow[d,"f'"]\\
S \& S' \arrow[l,"p"]
\end{tikzcd}}
cartésien \unsure{dans|ds} Cat. Si $f$ $W$-asphérique, $p$ $W$-lisse, alors $f'$ est $W$-asph.

\deleted{\unsure{Dém.} Dire que $f$ est asphérique signifie ici que $f\colon X\to Y$ est dans $W$ \unsure{d'apr.}}

Dém. \deleted{En \unsure{utilisant} la propriété a) de prop. 1\deleted{, \unsure{et se}} pour \unsure{prendre} la notion de $W$-asphéricité, et se rappelant que les isom. locaux sont lisses et par suite si $Y''\to Y'$ est un \unsure{iso|iso.} local, $Y''$ est lisse sur $Y$, on est \unsure{ra}} Il faut prouver $X'\times_{S'}S'_{/s'}$ est $W$-asphérique ($s'\in\mathrm{Ob}\, S'$). Soit $y=p(y')$, on a
\diagram{\begin{tikzcd}[ampersand replacement=\&,column sep=small,row sep=small]
S \& S' \arrow[l,"p"'] \& {}\\
{} \& \deleted{\unsure{S}}S_{/s}\arrow[ul] \& S'_{/s'}\arrow[l,"p_0"]\arrow[ul],
\end{tikzcd}}
avec $p_0$ \inserted{$W$}-lisse \inserted{(\unsure{limites|lemme} \unsure{b.ses})} et à fibres $W$-asphériques, donc une $W$-équiv. universelle (prop. 5). Comme l'image inverse de $X$ sur $X_{/y}$ est $W$-asph. (hyp. sur $f$), il

% archive sheet number: 82
\scanpage{83}
\mspage{32}

\diagram{\begin{tikzcd}[ampersand replacement=\&,column sep=small]
{} \& X_{/s}\arrow[d] \& X'_{/s'}\arrow[l,"W"']\arrow[d,no head]\\
(\unsure{\alpha|a}) \& \unsure{S|\mathscr S}_{/s} \& S'_{/s'}\arrow[l,"W"]
\end{tikzcd}}
que dans le carré cartésien ci-contre les flèches \unsure{horizon}\deleted{\unsure{tales}} sont dans $W$, et la première flèche verticale l'est par hypothèse sur $f$, \unsure{donc} la seconde aussi, donc $X'_{/y'}$ $W$-asphérique.

\deleted{\unsure{Prop. 7}} \uline{Corollaire} (généralise prop. 5, obtenue en faisant $Y=S$). Soit $p\colon S'\to S$ morphisme $W$-lisse dans Cat. Alors le foncteur de \unsure{cht.|chgt.} de base
\[
\mathrm{Cat}_{/S}\longrightarrow\mathrm{Cat}_{/S'}
\]
est compatible avec les localiseurs $W_S^{\unsure{g|9}}, W_{S'}^g$, i.e. transforme $f\colon X\to Y$ qui est dans $W$ \uline{localement} sur $S$, en $f'\colon X'\mathbin{\unsure{\text{---}|\to}}Y'$ \deleted{\unsure{dans}} $\in W$ localement sur $S'$.

Cette fois, on regarde le diagramme similaire à $(\alpha)$
\diagram{\begin{tikzcd}[ampersand replacement=\&,row sep=small,column sep=small]
X_{/s}\arrow[dr,"f_{/s}"']\arrow[dd,no head] \& {} \& X'_{/s'}\arrow[ll,"W"']\arrow[dr,"f'_{/s'}"]\arrow[dd] \& {}\\
{} \& Y_{/s}\arrow[dl,no head] \& {} \& Y'_{/s'}\arrow[ll,"W"',crossing over]\arrow[dl]\\
S_{/s} \& {} \& S'_{/s'}\arrow[ll,"W"'], \& {}
\end{tikzcd}}
où les flèches horizontales sont dans $W$, donc $f_{/s}\in W$ (hyp.) implique $f'_{/s'}\in W$.

% archive sheet number: 83
\scanpage{84}
\mspage{33}

Dualement, si $p$ est $W$-propre, le foncteur image inverse est compatible avec les localiseurs $W_S^d, W_{S'}^d$.

Si $f\colon X\to S$ est $W$-propre, on va définir un \deleted{\unsure{objet|foncteur}} \inserted{foncteur} canonique
\[
\begin{aligned}
\deleted{\unsure{\mathrm{hot}_{W, *}(\mathrm{Hot}_{W})=\mathrm{hot}^{X}_{*}}}
\inserted{\mathrm{hot}^{X/S,W}_{*}}\colon S&\longrightarrow\mathrm{Hot}_W\overset{\mathrm{def}}{=}(\mathrm{Cat})W^{-1}\\
\text{donné sur les objets par}\qquad s&\longmapsto\mathrm{hot}_W(X_s).
\end{aligned}
\]
Il faut définir, \deleted{\unsure{pour}} une flèche $u$ dans $S$, $u\colon s\longrightarrow s'$, une flèche
\[
\mathrm{hot}_W(X_s)\mathbin{\unsure{\text{---}|\to}}\mathrm{hot}_W(X_{s'})
\]
dans $\mathrm{Hot}_W$ -- flèche que nous pouvons noter $\deleted{\unsure{(}}\mathrm{hot}_{\deleted{\unsure{W}}*}^{X/S,W}\deleted{\unsure{)}}(u)$. On la définit par le diagramme
\diagram{\begin{tikzcd}[ampersand replacement=\&,column sep=large]
X_s\arrow[r,dashed]\arrow[d,"i_s^X"] \& X_{s'}\arrow[d,"i_{s'}^X"]\\
X_{/s}\arrow[r,"X_{/u}"] \& X_{/s'},
\end{tikzcd}}
où les flèches verticales sont $W$-\unsure{coasph.|coasph} donc dans $W$. Ce sont \inserted{donc} des iso dans $\mathrm{Hot}_W$, d'où l'unique flèche en pointillés dans $\mathrm{Hot}_W$, rendant le carré commutatif. Sur cette définition il est immédiat que
\[
u\longmapsto\mathrm{hot}_W(u)
\]
est compatible avec composition, et unités. Ce foncteur est \inserted{donc} défini \deleted{\unsure{sous}} \deleted{\unsure{des}} supposant

% archive sheet number: 84
\scanpage{85}
\mspage{\unsure{34}}

seulement que les
\[
X_y\mathbin{\unsure{\text{---}|\to}}X_{/y}
\]
sont dans $W$, ce qui est nettement moins fort que propre. Si on suppose que cette hypothèse est
% slanted
\margin{On dira que $f$ est $W$-\uline{prépropre}}
satisfaite ``universellement'' (i.e. reste satisfaite après tout chgt de base $\unsure{S|\mathscr S}'\mathbin{\text{---}}S$), alors on a donc, pour tout $S'$ dans $\mathrm{Cat}_{/S}$, un \deleted{\unsure{fonction}} foncteur
\[
\mathrm{hot}^{X'/S',W}_{\deleted{\unsure{W}}*}\colon S'\longrightarrow\mathrm{Hot}_W,
\]
et on vérifie \deleted{\unsure{tout}} de suite que ce n'est autre que le composé dans
\diagram{\begin{tikzcd}[ampersand replacement=\&,row sep=large]
{} \& S'\arrow[rr,"p"]\arrow[dr,"{\mathrm{hot}_W^{X'/S'}}"'] \& {} \& S\arrow[dl,"{\mathrm{hot}_W^{X/S}}"]\\
(\unsure{\alpha|a|*}) \& {} \& \mathrm{Hot}_W \& {.}
\end{tikzcd}}
On a des observations duales si $X$ est lisse sur $S$, plus généralement si les inclusions
\[
X_s\mathbin{\text{---}}{}_{s\backslash}X
\]
sont dans $W$. On trouve alors un contrafoncteur
\[
\left\{\begin{aligned}
S^\circ&\longrightarrow \mathrm{Hot}_W\\
s&\longmapsto\mathrm{hot}_W^{\deleted{\unsure{*}}}(X_s).\\
u&\mathrel{\unsure{\vdash|\mapsto}}(\mathrm{hot}^{\deleted{\unsure{X/S,W}}}_{\unsure{X/S},W})^*(u)\quad\text{ou}\quad\mathrm{hot}_{\unsure{,|f}X/S}^{*}
\end{aligned}\right.
\]

% archive sheet number: 85
\scanpage{86}
\mspage{35}

et une compatibilité similaire à $(\unsure{\alpha|a|*})$, si l'hyp. est satisfaite \uline{universellement} \unsure{sur} $S$, \unsure{p. ex.|par ex.} $f$ $W$-lisse.

\uline{Proposition 7} Supposons $f\colon X\to\unsure{S|\mathscr S}$ à la fois $W$-propre et $W$-lisse, plus gén. que les morphismes
\[
X_s\mathbin{\unsure{\text{---}|\to}}X_{/s},\qquad X_s\mathbin{\text{---}}{}_{s\backslash}X\qquad(s\in S)^{\unsure{\prime|\prime\prime}}
\]
soient dans $W$, et que cette hypothèse reste vérifiée après tt chgt de base. Alors les foncteurs correspondants
\[
\begin{aligned}
\lambda\deleted{^{\unsure{X|f}}_{\unsure{*}}}&=\mathrm{hot}^{X/S,W}_{\deleted{\unsure{W}}*}\colon S\longrightarrow\mathrm{Hot}_W\\
\lambda^\circ&=\mathrm{hot}_{X/S,W}^{*}\colon S^\circ\mathbin{\text{---}}\mathrm{Hot}_W
\end{aligned}
\]
transforment \unsure{toute} flèche \inserted{$u$ de $S$} en iso. de $\deleted{\unsure{H}}\mathrm{Hot}_W$, et on a
\[
\lambda^\circ(u)=\lambda(u)^{-1}\qquad\text{dans }\mathrm{Hot}_W,
\]
i.e. $\lambda(u)$ et $\lambda^\circ(u)$ sont inverses l'un de l'autre.

\uline{Dém} La donnée de $u\colon\inserted{s_0\to s_1}$ équivaut à celle de
\[
S'=\mathbb{\Delta}^1\longrightarrow S,
\]
d'où $X'=X\times_S S'$. sur $S'=\mathbb{\Delta}^1$, avec les deux fibres
\[
X'_0\simeq X_{s_0},\qquad X'_1\simeq X_{s_1}.
\]
L'hyp. assurant l'existence de $\lambda$ s'exprime par la condition que
\[
X'_1\longrightarrow X'_{/1}\simeq X'
\]

% archive sheet number: 86
\scanpage{87}
\mspage{36}

est dans $W$, de sorte que le diagramme
\diagram{\begin{tikzcd}[ampersand replacement=\&]
X'_0\arrow[r,"i_0","{\deleted{\unsure{W}}}"'] \& X' \& X'_1\arrow[l,"i_1"',"W"]
\end{tikzcd}}
définit une flèche \inserted{$\lambda(u)\colon$} $X'_0\to X'_1$\unsure{,} dans $\mathrm{Hot}_W$, savoir $\lambda(u)=i_1^{-1}\,i_0$. L'hyp. qui assure l'existence de $\lambda^\circ$ s'exprime par la condition symétrique $i_0\in W$, qui permet de définir la flèche $\lambda^\circ(u)\colon X'_1\mathbin{\unsure{\text{---}|\to}}X'_0$ par $\lambda^\circ(u)=i_0^{-1}i_1$.

Quand les deux conditions sont satisfaites simultanément, on voit donc que $\lambda(u),\lambda^\circ(u)$ \unsure{sont} inverses \unsure{l'un|l'une} de l'autre.

\uline{Remarques} Pour définir $\lambda$ ``universellement'', il n'était pas indispensable de considérer les $X_{/s}$ et les inclusions $X_s\mathbin{\text{---}}X_{/s}$, au lieu de cela il suffisait de regarder les chgts. de base $S'=\mathbb{\Delta}^1\xrightarrow{u}S$, et de faire l'hyp. que, pour tout tel chg. de base, l'inclusion
\[
X'_1\xrightarrow{i_1}X'=X_{(\mathbb{\Delta}^{1},u)}
\]
est dans $W$. Cette hypothèse est automatiquement stable par chg. de base arbitraire. (La $W$-propreté de $f\colon X\to S$ signifie que $X'_1\mathbin{\text{---}}X'$ est non seulement dans $W$, mais $\unsure{=}$ $W$-coasphérique.) Dualement, pour la définition ``universelle'' de $\lambda^\circ$, il \unsure{suffirait|suffisait} que pour tt chgt de base $\mathbb{\Delta}^1\to S$ comme dessus, l'inclusion $X'_0\xrightarrow{i_0}X'$ est dans $W$.

% archive sheet number: 87
\scanpage{88}
\mspage{37}

\deleted{C'est \unsure{ça|ce} \unsure{que|qui} \unsure{me semble}} La notion la plus faible raisonnable de $W$-fibré, me semble être celle où on suppose les deux choses à la fois.

Les hyp. qu'on a faites jusqu'à présent, à l'exclusion de la forme forte \deleted{\unsure{$f$}} de $W$-propreté ou de $W$-lissité, reviennent toujours à supposer que le \uline{\unsure{changmt|changement} de base par} $f\colon X\mathbin{\unsure{\text{---}|\to}}S$ transforme certaines flèches dans $\mathrm{Cat}_{/S}$ en $\in W$\,: Pour définir $\lambda$, ces flèches sont
\[
\text{soit}\quad\{s\}\to S_{/s},\qquad\text{soit}\quad\{1\}\mathbin{\text{---}}\mathbb{\Delta}^{1}\;(\to S),
\]
pour définir $\lambda^\circ$, ce sont
\[
\text{soit}\quad\{s\}\mathbin{\text{---}}{}_{s\backslash}S\deleted{\unsure{_{s}}},\qquad\text{soit}\quad\{0\}\mathbin{\text{---}}\mathbb{\Delta}^1\;(\mathbin{\text{---}}S),
\]
le ``soit'' suivant la définition qu'on a adoptée. Dans ces exemples, il s'agit toujours de l'inclusion de l'objet final d'une catégorie $S'$, dans \unsure{la dite}, ou inversement de l'objet initial.

Cela montre d'ailleurs (si on exige l'universalité), que \unsure{la|ladite} première hypothèse (\unsure{assurant|concernant} $\lambda$) est \uline{exactement} celle (\inserted{cas de l'}inclusion d'un objet final \unsure{$\}s'$|$)s'$} -- puisque pour un tel objet, on a $S'_{/s'}\simeq S'$), tandis que la deuxième hypothèse exige moins, seulement pour le cas de $\mathbb{\Delta}^1$. Et dualement pour \deleted{\unsure{$\lambda$}} \unsure{$\lambda'$|$\lambda^\circ$}.

Mais il convient de noter que pratiquement, la façon la plus commode, et de loin,

% archive sheet number: 88
\scanpage{89}
\mspage{38}

pour vérifier de telles propriétés, c'est soit par la $W$-propreté de $f$, soit par la $W$-lissité -- ce qui revient à vérifier que certaines catégories bien précisées sont $W$-asphériques, sans avoir à se soucier de chgts de base et de vérifier que des flèches construites en termes de chgt. de base arbitraires sont dans $W$. Et l'emploi de la $W$-propreté et de la $W$-lissité est d'autant plus sympa, que ces notions sont stables par composition, alors que \deleted{\unsure{toutes}} les notions \unsure{dérivées}, ainsi que les notions de fibration à la clef, ont une fâcheuse tendance à ne \uline{pas} être stables par composition.

On pourrait appeler pseudo-\inserted{$W$-}propres resp. quasi-\inserted{$W$-}propres (on \unsure{verra} mieux que pseudo) ce qu'on vient de dégager) -- j'entends des \unsure{ricanements}! Et si on veut les deux à la fois -- $\lambda$ et $\lambda^\circ$ \unsure{associés} comme on a dit? Je vois le $\unsure{=}$ principe de description avec \inserted{des} $S''\to S'$ sur $S$, d'exiger que $X''\mathbin{\unsure{\text{---}|\to}}X'$ soit transformé en des $\in W$.

\unsure{Avec|Donc|D'où} les choix suivants qui ont l'air sympa.

% archive sheet number: 89
\scanpage{90}
\mspage{39}

$a_1)$\quad $Y'$ a un objet final, $Y''=e$. La \unsure{définition} revient \deleted{\unsure{donc}} à la condition que les
\[
X_{\unsure{s|y}}\mathbin{\unsure{\text{---}|\to}}X_{/s},\qquad X_{/\deleted{\unsure{y}}s'}\longrightarrow X_{/\deleted{\unsure{y}}s}
\]
(pour $s\in\mathrm{Ob}\,S$, où $u\colon s'\mathbin{\text{---}}s$ dans $S$) soient dans $W$, et que ça reste comme ça après tout chgt. de base $S'\to S$.
% slanted
\margin{NB $a_1),a_2)$ sont \unsure{notés|notées} (\unsure{pour a)|par ex.} p. 4, mais \unsure{$a_0$|$a_2$} a été oublié p. 4)}
$a_2)$\quad $Y'$ a un objet initial, $Y''=e$, c'est la condition duale à coups de
\[
X_s\longrightarrow{}_{s\backslash}X,\qquad{}_{s\backslash}X\longrightarrow{}_{s'\backslash}X.
\]
$\deleted{\unsure{a}}\unsure{b_0|a_0})$\quad $Y'=\mathbb{\Delta}^1$, $Y''=e$.

$\deleted{\unsure{a}}b)$\quad $Y'$ est contractile, $Y''$ aussi

(ou au choix, $Y'$ contractil, $Y''=e$) donc la condition revient à ceci\,: pour tout $Y'$ \inserted{sur $Y$} \uline{contractile} (notion indépendante de $W$), c'est ça son principal avantage) les inclusions des fibres de $X'=X\times_Y Y'$ sur $Y'$ dans $X'$ \deleted{sont} $\in W$.

Je sens que même pour des $W$ générales, $\unsure{b})$ \unsure{reste|sort|est} du domaine du vérifiable, tandis qu'on \unsure{mollit} si on passe à l'extrême en généralité dans ce sens\,:

$b')$\quad Comme $b)$, \deleted{\unsure{mais}} \deleted{\unsure{en}} en supposant seulement $Y'$ $W$-asphérique.
% slanted
\margin{C'est \deleted{\unsure{le|ess.}} le $a')$ de p. 4}
Enfin (la condition a de la p. 4)

$c)$\quad Si $Y''\to Y'$ est $\in W$ \deleted{\unsure{resp.}} sans plus, $X''\to X'$ aussi,

% archive sheet number: 90
\scanpage{91}
\mspage{40}

La seule de ces notions qui ait la moindre envie d'être stable par composition, c'est la plus forte $c)$. Par contre, la notion de flèche $W$-parfaite, \unsure{(cf} déf. 3) est stable spontanément par composition\,:

\uline{Définition 3} On dit que $f\colon X\mathbin{\unsure{\text{---}|\to}}S$ est $W$-\unsure{parfaite|parfait}, si elle est à la fois $W$-propre et $W$-lisse.

Dans le cas $W=W_\infty$, toutes les variantes de la \unsure{page|prop.} \unsure{précédentes|précédente}, à l'exclusion seulement de $a_0$, sont équivalentes. Ce serait intéressant de dégager des conditions \inserted{minimales} sur $W$ pour assurer qu'il en soit ainsi. (Et \deleted{\unsure{de}} se \unsure{pose|poser} la question si dans le cas de $W_\infty$, $a_0$ suffit pour impliquer le reste -- je ne m'y attends guère \ldots)

Dans le cas général, on a le diagramme d'implications (et encore je me suis limité)
\diagram{\begin{tikzcd}[ampersand replacement=\&,row sep=small,column sep=small,execute at end picture={
\draw[thin] ([xshift=8pt,yshift=2pt]bb.east) .. controls +(-0.5,0) and +(0.9,0.1) .. ([yshift=3pt]parf.north east) .. controls +(-0.55,0.35) and +(-0.55,0.35) .. ([xshift=-5pt]parf.west) .. controls +(-0.1,-0.45) and +(-0.3,0) .. ([yshift=-4pt]parf.south east) .. controls +(0.6,-0.1) and +(-0.25,0.1) .. ([xshift=-2pt,yshift=-5pt]bb.south west) .. controls +(0.55,-0.28) and +(0.25,-0.2) .. ([xshift=9pt,yshift=-1pt]bb.east);}]
{} \& {} \& c\arrow[d,Rightarrow] \& {} \& \text{\scriptsize $W$-fibration ultra-forte}\\
|[alias=parf]|\text{\scriptsize $W$-parfait}\arrow[rrd,Rightarrow] \& {} \& b'\arrow[d,Rightarrow] \& {} \& \text{\scriptsize \textquotedbl\quad $W$-forte}\\
{} \& {} \& |[alias=bb]|b\arrow[dl,Rightarrow]\arrow[dr,Rightarrow] \& {} \& \text{\scriptsize \textquotedbl\quad hot-forte}\\
\text{\scriptsize $W$-fibr. gauche} \& a_1\arrow[dr,Rightarrow] \& {} \& a_2\arrow[dl,Rightarrow] \& \text{\scriptsize $W$-fibr. droite}\\
{} \& {} \& a_0 \& {} \& {}\\
{} \& {} \& \text{\scriptsize $W$-pseudo fibration} \& {} \& {}
\end{tikzcd}}
\margin{(Il semblerait que ce soit ça, finalement, la ``bonne notion'' de $W$-fibration, \uline{équivalente} à $a_1$ et $a_2$, lesquelles seraient équivalentes entre elles. Cf à propos p. \unsure{48}.}

% archive sheet number: 91
\scanpage{92}
\mspage{41}

\uline{Proposition} 8\,: Soit $f\colon X\longrightarrow S$ dans Cat, \deleted{$f$ $W$-propre} \textsuperscript{*}
% slanted
\margin{\deleted{Il suffit $f$ $W$-pré-propre i.e. les $X_{\unsure{s|y}}\to X_{/s}$ dans $W$, ou bien il ne \unsure{reste|sera} de \unsure{pr.} \unsure{que|de}}}
Les conditions suivantes sont équivalentes

a) $f$ est une $W$-fibration \uline{gauche}, i.e. la condition $a_1)$ de la page 39 est vérifiée (qui implique

b) \inserted{$f$ est $W$-prépropre, et} $\forall\,s_0\to s_1$ dans $S$, $X_{/s_0}\mathbin{\unsure{\text{---}|\to}}X_{/s_1}$ est dans $W$.

c) $f$ \inserted{est $W$-prépropre, et} est une $W$-pseudofibration, i.e. la condition $a_0$ de la page 39 est vérifiée.
% slanted
\margin{\uline{Cor} Quand $f$ est Cat-\unsure{cofibrant|cofibrante}, alors ces conditions équivalent aussi à ceci\,: $\forall\,u\unsure{\in|:}s_0\to s_1$ dans $S$, $u_{\unsure{!|*}}\colon X_{s_0}\mathbin{\text{---}}X_{s_1}$ est dans $W$.}
\uline{NB} Ces conditions sont satisfaites si $f$ est en $\unsure{=}$ temps $W$-pré-lisse, cf prop. 7.

\uline{Dém} On a $a)\Longrightarrow b)$ trivialement (faire $\unsure{Y'|Y''}=Y_{/s}$, $Y''=Y_{/s'}$). $b\Longleftrightarrow c$ par le raisonnement de la prop. 7 car en termes de $X'=X_u=X\times_S(\mathbb{\Delta}^1,u)$, la condition b) pour $u\colon s_0\mathbin{\text{---}}s_1$ \unsure{donné|donnée} s'interprète \deleted{que} \inserted{dans} le diagramme
\diagram{\begin{tikzcd}[ampersand replacement=\&,row sep=small,column sep=small,execute at end picture={
\draw[thin] ([xshift=-1pt]note.west) -- ([yshift=4pt]za.north) -- (note.east) -- ([yshift=2pt]zb.north) -- (zca.west) -- (zr.east) -- ([xshift=-5pt,yshift=-8pt]zc.south) -- (zrc.east) -- (zcb.west);\draw[thin] (zc.south east) -- ([xshift=-3pt,yshift=4pt]zx.north) -- ([xshift=-12pt]zx.north west) .. controls +(0.8,0.15) and +(-0.4,0.1) .. (zx.north east) -- (zcb.south);} ]
{} \& |[alias=za]|\deleted{X'_0}\arrow[r,"{\deleted{i_0}}"]\arrow[d,no head,"{\deleted{\simeq}}"] \& |[alias=zb]|\deleted{X'} \& |[alias=zr]|\deleted{X'_1}\arrow[l,"{\deleted{i_1\in W}}"']\arrow[d,no head,"{\deleted{\simeq}}"]\\
|[alias=note]|\substack{\deleted{i_1\in W}\ \text{\deleted{car $X'$ propre}}\\ \text{\deleted{sur $S'=\mathbb{\Delta}^1$}}} \& |[alias=zca]|\deleted{X_{s_0}}\arrow[d,"{\deleted{\in W}}"] \& {} \& |[alias=zrc]|\deleted{X_{s_0}}\arrow[d,"{\deleted{\in W}}"]\\
{} \& |[alias=zc]|\deleted{X_{/s_0}}\arrow[rr]\arrow[dr,no head] \& {} \& |[alias=zcb]|\deleted{X_{/s_1}}\arrow[dl]\deleted{.}\\
{} \& {} \& |[alias=zx]|\deleted{\unsure{X}} \& {}
\end{tikzcd}}
\unsure{comm.|comme.}
\diagram{\begin{tikzcd}[ampersand replacement=\&,column sep=small,row sep=small]
X'_0\arrow[dr,"{i'_0\in W}","W"']\arrow[dd,"{\sim}"] \& {} \& X'_1\arrow[dr,"{i'_1\in W}","{\unsure{W}}"']\arrow[dd,no head,"{\sim}"'] \& {}\\
{} \& \substack{X'_{/0}\\(\simeq X'_0)}\arrow[rr,"{\varphi'}",crossing over]\arrow[dd,"p_0"] \& {} \& X'_{/1}\;(\simeq X')\arrow[dd,"p_1"]\\
X_{s_0}\arrow[dr,"{i_0\in W}","W"'] \& {} \& X_{s_1}\arrow[dr,"{i_1\in W}"] \& {}\\
{} \& X_{/s_0}\arrow[rr,"{\varphi}"] \& {} \& X_{/s_1},
\end{tikzcd}}
où $i_0,i_1,i'_0,i'_1$ sont dans $W$ à cause de l'hypothèse de propreté (stable par chgt. de base) en fait il suffisait \deleted{l'hyp.} cette condition de $W$-pré-propreté. Donc $p_0,p_1$ sont $\in W$, donc $\varphi\in W$ ssi $\varphi'\in W$, \deleted{ce qui} \inserted{or $\varphi'\in W$} est la condition $a_0$.

% archive sheet number: 92
\scanpage{93}
\mspage{42}

Prouvons enfin que $b,c$ impliquent $a)$. Comme $c)$ \inserted{(\unsure{comme|ou} aussi $b$)} est stable par chgt. de base, il suffit de montrer ceci\,: Si $S$ a un objet final $\deleted{\unsure{s}}e$, alors pour tout $s\in\mathrm{Ob}\,S$, l'inclusion $X_s\mathrel{\overset{i}{\unsure{\text{---}|\to}}}X$ est dans $W$. Or on a le diagramme commutatif
\diagram{\begin{tikzcd}[ampersand replacement=\&,row sep=large,column sep=large,execute at end picture={\draw[thin] (old.west) -- (old.north east) -- (old.south west) -- ([xshift=3pt]old.north east) -- (old.south east) -- ([xshift=8pt]old.north east) -- ([xshift=4pt]old.south east) -- ([xshift=14pt]old.north east) -- ([xshift=9pt]old.south east) -- ([xshift=-10pt]target.north west) -- ([xshift=-16pt]target.west) -- ([xshift=-4pt]target.north west);}]
X_s\arrow[r,"{i_s\in W}"]\arrow[dr,"i"] \& X_{/s}\arrow[d,"{u\in W\ \text{par }c)}"]\\
|[alias=old]|\deleted{\unsure{X_e}}\arrow[r,"{\deleted{\unsure{i_e\in W}}}"] \& |[alias=target]|X_{/e}=X,
\end{tikzcd}}
qui montre bien $i=ui_s\in W$.

\deleted{Mais} Dualement, si $f$ est $W$-prélisse, alors \deleted{\unsure{$f$}} \deleted{\unsure{c'est}} la condition $a_0$ ($f$ une $W$-pseudofibration) est équivalente à $a_2$ ($f$ une $W$-fibration droite), et signifie aussi que les
\[
{}_{s_{\unsure{1|0}}\backslash}X\mathbin{\text{---}}{}_{s_0\backslash}X
\]
sont dans $W$. Mais la notion de ``$W$-fibration'' sur laquelle on tombe (si on ne se contente du misérable $a_0$!) est différente dans les deux cas duaux l'un de l'autre. On aimerait dans les deux cas trouver des $W$-fibrations ``hot-fortes'' (condition $b$), tout au moins en mettant du piquant (s'il le faut) -- i.e. pour $f$ $W$-parfait. Ça revient donc à ceci

\uline{Question}\,: Soit $X\xrightarrow{f}S$, avec $f$ $W$-parfait, $S$ contractile, $s\in\mathrm{Ob}\,S$, L'inclusion $X_s\to X$ est-elle dans $W$\,?

% archive sheet number: 93
\scanpage{94}
\mspage{43}

On va supposer \inserted{seulement} que $f$ est $W$-lisse et une $W$-fibration gauche \inserted{(cf prop. \unsure{7|8})}. On va se ramener au cas où $X$ est même fibrant sur $S$, en introduisant $\unsure{\Phi|\mathscr L}_S(X)$ et
\[
X\mathrel{\overset{i_X}{\unsure{\text{---}|\to}}}\Phi_S(X)
\]
L'ennui, c'est que j'ai fait la théorie de $\mathrm{HOT}_W$ jusqu'à maintenant sous l'hyp. (qui s'avère \unsure{idiote}) que $W\supset W_\infty$ -- je serai donc obligé de revoir rapidement sous quelles hyp. plus vraisemblables sur $W$ les résultats familiers sont valables. \deleted{J'aimerais \unsure{voir}} Par exemple j'aimerais que $i_X\in W$, de sorte qu'on ait
% slanted
\margin{en fait, $\unsure{j|i}_X\in(W_\infty)_S^d\subset W_S^d\unsure{,}\subset W_{\unsure{S|\}}}$}
\diagram{\begin{tikzcd}[ampersand replacement=\&,column sep=huge,row sep=large]
X_s\arrow[r,"{\substack{\deleted{\unsure{i_X}}\\i_s\in W}}","{\text{car $X$ $W$-lisse}} "']\arrow[d,"j_s^X"'] \& \Phi_S(X)_s={}_{s\backslash}X\arrow[d,"{j_s^{\Phi_S(X)}}"]\\
X\arrow[r,"{i_X\in W}"] \& \Phi_S(X)
\end{tikzcd}}
on voit donc que pour prouver que l'inclusion $j_s^X$ est $\in W$, il suffit de prouver la $\unsure{=}$ chose pour $\Phi_S(X)$, qui est fibrant \inserted{(pas seulement \unsure{$W$-lisse})}. De plus, l'hypothèse pseudo-fibration sur $X$ passe sûrement à $\Phi_S(X)$ \unsure{a}. (je réserve les deux vérifications à plus tard).
% slanted
\margin{OK}
On est donc ramené au cas \unsure{cas|où} de
\[
X\in\mathrm{Ob}\,\underline{\mathrm{Fib}}_0\, S,\qquad S\text{ contractile},
\]
à prouver que l'inclusion
\[
\deleted{\unsure{s|x}}X_s\longrightarrow X
\]

% archive sheet number: 94
\scanpage{95}
\mspage{44}

est dans $W$. Je vais prouver ceci\,:

\uline{Lemme} \unsure{Soient} $S$ dans Cat, $X\in\mathrm{Ob}\,\underline{\mathrm{Fib}}_0\,S$, et considérons, pour un $S'$ dans Cat,
\[
\begin{aligned}
f&\mathrel{\unsure{\vdash|\mapsto}}f^*(X/S)=X\times_S(S',f)\\
\underline{\mathrm{Hom}}(S',S)^\circ&\longrightarrow\underline{\mathrm{Fib}}_0\,S'
\end{aligned}
\]
(Ce foncteur est défini grâce au fait que $X$ est Cat-fibrant sur $S$, et peut s'interpréter donc, \deleted{\unsure{à une}} \inserted{à $S$-}équivalence près, en termes d'un foncteur $S^\circ\mathbin{\unsure{\text{---}|\to}}\mathrm{Cat}$ \ldots) Les flèches de $\underline{\mathrm{Hom}}(S',S)$ sont transformées en flèches $\in W^d_{S'}$. \deleted{\unsure{donc}}
% slanted
\margin{\uline{NB} $\unsure{\Phi|\mathscr L}_S(X)$ était \unsure{même|justement} déjà \unsure{scindé|scindée}, \unsure{définie|défini} par le foncteur $s\mapsto{}_{s\backslash}X_{\deleted{\unsure{s}}}$}
En effet, en termes de foncteurs
\[
S'{}^\circ\longrightarrow\mathrm{Cat}
\]
si $X$ est défini par
\[
S^\circ\xrightarrow{\unsure{\mathscr X|\mathcal X}}\mathrm{Cat},
\]
alors $f^*(X)$ est défini par $\mathscr X\circ f$, et la flèche
\[
\mathscr X\circ f\mathbin{\text{---}}\mathscr X\circ g
\]
définie par $u\colon f\mathbin{\text{---}}g$ n'est autre que $\mathscr X*u$. Sa valeur en $s'\in\mathrm{Ob}\,S'$ est obtenue en prenant
\[
u(s')\colon f(s')\longrightarrow g(s')
\]
d'où par $\mathscr X$
\[
\mathscr X(u(s'))\colon
\underset{\deleted{\unsure{g}}(\mathscr X\circ g)(s')}{\tikz[baseline=(term.base)]{\node[inner sep=0pt](term){$\mathscr X(g(s'))$};\draw ([xshift=-1pt]term.south west) -- ([xshift=-1pt,yshift=-2pt]term.south west) -- ([xshift=1pt,yshift=-2pt]term.south east) -- ([xshift=1pt]term.south east);}}
\longrightarrow
\underset{(\mathscr X\circ f)(s')}{\tikz[baseline=(term.base)]{\node[inner sep=0pt](term){$\mathscr X(f(s'\unsure{)|,})$};\draw ([xshift=-1pt]term.south west) -- ([xshift=-1pt,yshift=-2pt]term.south west) -- ([xshift=1pt,yshift=-2pt]term.south east) -- ([xshift=1pt]term.south east);}}
\]
% slanted
\margin{à dégager en énoncé à part}
Mais l'hyp. \deleted{que} $X\in\underline{\mathrm{Fib}}_0(S)$ signifie que les $\mathscr X(w)$, $w\in\mathrm{Fl}(S)$, sont \unsure{ts|tous} dans $W$. \unsure{Donc|D'où}

% archive sheet number: 95
\scanpage{96}
\mspage{45}

$\unsure{\mathscr X|\mathcal X}\unsure{*u}$ est dans $W$ fibre par fibre, donc dans $W_{S'}^d$.

\uline{Corollaire 1} Le foncteur du lemme, composé avec
\[
\underline{\mathrm{Fib}}_0\,S'\longrightarrow\mathrm{HOT}_W^{\unsure{\mathrm{lc}|\mathrm{le}}}(S')
\overset{\mathrm{def}}{=}(\underline{\mathrm{Fib}}_0\,S')(W_{S'}^d)^{-1},
\]
se factorise par le groupoïde fondamental de $\underline{\mathrm{Hom}}$\,:
\diagram{\begin{tikzcd}[ampersand replacement=\&,row sep=large,column sep=large]
\underline{\mathrm{Hom}}(S',S)^\circ\arrow[r]\arrow[d] \& \underline{\mathrm{Fib}}_0\,S'\arrow[d]\\
\Pi\bigl(\underline{\mathrm{Hom}}(S',S)\bigr)^\circ\arrow[r] \& \mathrm{HOT}_W^{\mathrm{lc}}(S')
\end{tikzcd}}
Si $f,g$ sont dans une $\unsure{=}$ composante connexe de $\underline{\mathrm{Hom}}(S,S')$, alors $f^*X$ et $g^{\unsure{\prime|*|\wedge}}X$ sont isomorphes \deleted{dans} (non canoniquement, en général) dans $\mathrm{HOT}_W^{\mathrm{lc}}(S')$.

\deleted{\uline{Corollaire} \deleted{\unsure{Soit}} \unsure{Supposons}} \inserted{Il est plus astucieux de ne pas fixer $X$,} et \deleted{\unsure{d'écrire|décrire}} considérer
\[
\underline{\mathrm{Hom}}(S',S)\longrightarrow\underline{\mathrm{Hom}}\bigl(\mathrm{HOT}_W^{\mathrm{lc}}(S),\mathrm{HOT}_W^{\mathrm{lc}}(S')\bigr)\tag*{\makebox[0pt][r]{\makebox[\displaywidth][l]{($\unsure{\alpha|a}$)}}}
\]
correspondant à
\[
\begin{aligned}
\underline{\mathrm{Hom}}(S',S)\times\mathrm{HOT}_W^{\mathrm{lc}}(S)&\longrightarrow\mathrm{HOT}_W^{\mathrm{lc}}(S')\\
(f,\unsure{\xi|\zeta})&\mathrel{\unsure{\vdash|\mapsto}}f^*(\xi).
\end{aligned}
\]
Le lemme dira que les flèches de $S'\mathbin{\unsure{\text{---}|\to}}S$ sont transformées par ($\alpha$) en isomorphismes, \unsure{non}, d'où
\[
\Pi\,\underline{\mathrm{Hom}}(S',S)\longrightarrow\underline{\mathrm{Hom}}\bigl(\mathrm{HOT}_W^{\mathrm{lc}}(S),\mathrm{HOT}_W^{\mathrm{lc}}(S')\bigr),
\]

% archive sheet number: 96
\scanpage{97}
\mspage{46}

donc si $f,g\colon S'\mathrel{\unsure{\substack{\longrightarrow\\[-.8ex]\rule{1.25em}{.35pt}}|\rightrightarrows}}S$ sont deux foncteurs homotopes, alors
\[ f^*\simeq g^*\colon \mathrm{HOT}^{\unsure{\mathrm{lc}|\mathrm{le}}}_W(S)\rightrightarrows\mathrm{HOT}^{\mathrm{lc}}_W(S') \]
Il s'ensuit aussitôt ceci

\uline{Corollaire 2} Si $f$ est un homotopisme, $f^*\colon\mathrm{HOT}^{\mathrm{lc}}_W(S)\unsure{\text{---}|\to}\mathrm{HOT}^{\mathrm{lc}}_W(S')$ est une équivalence de catégories.

En effet, prenant $g\colon S\longrightarrow S'$ \unsure{telle|tel} que
\[ fg\sim\mathrm{id}_S,\qquad gf\sim\mathrm{id}_{S'}, \]
on trouve que
\[ \begin{aligned}\deleted{\unsure{f}}g^*f^*&\simeq(\mathrm{id}_S)^*=\mathrm{id}_{\mathrm{HOT}^{\mathrm{lc}}_W(S)}\\
f^*g^{\unsure{\prime|*|\wedge}}&\simeq(\mathrm{id}_{S'})^{\unsure{\prime|*}}=\mathrm{id}_{\mathrm{HOT}^{\mathrm{lc}}_W(S')},\end{aligned} \]
donc on trouve que (quitte à ajuster ces \unsure{isoms|isom} fonctoriels) $f^{\unsure{\cdot|*|\prime}}$ et $g^{\unsure{\prime|*}}$ sont des équivalences quasi-inv. l'une de l'autre.

J'ai l'impression que tout ça ne doit pas utiliser plus que les axiomes \unsure{standart|standard} sur $W$ comme localiseur fondamental, tout au plus $W(2)$ au lieu de la forme faible $W(2\mathrm{bis})$.
\margin{OK} % slanted
Revenons à $X$ lisse sur $S$ contractible. Le corollaire 2 indique que
\[ \begin{aligned}\mathrm{Hot}_W&\longrightarrow\mathrm{HOT}^{\mathrm{lc}}_W(S)\\
\unsure{\xi|\zeta}&\text{---}p_S^*(\xi),\end{aligned} \]
où $p_S\colon S\longrightarrow e$, est une équivalence. Donc dans la cat. $\mathrm{HOT}^{\mathrm{lc}}_W(S)$, $X$ est isomorphe

% archive sheet number: 97
